\SourceLeaf{317}%
Endlich findet sich
\[
\begin{aligned}
\log\lambda&=2.7315487\\
\log l&=2.7315438\\
l&=538''\,9442=0^0\,8'\,58''\,9442.
\end{aligned}
\]

\indent Die Bequemlichkeit dieses Verfahrens wird allerdings erst dann in ihrer vollen Grösse fühlbar, wenn man sich die Hülfen des kleinen Mechanismus bei Handhabung derartiger Methoden zu eigen gemacht hat, wozu eine Anweisung hier nicht an ihrem Platze sein würde. Ich begnüge mich hier nur anzudeuten, dass, was in obigem Beispiele wie eine viermalige Rechnung erscheint, nicht in der Form von vier getrennten Rechnungen, sondern wie eine einzige geschrieben werden soll, indem man bei jeder neuen Überarbeitung nur die letzten Zifern ergänzt oder verbessert. Jedenfalls braucht man immer nur die letzte Rechnung aufzubewahren, und gerade darin besteht ein grosser Vortheil, zumal bei Messungen von bedeutendem Umfange, dass man dann den ganzen wesentlichen Kern der Berechnung für alle Dreiecksseiten im möglich kleinsten Raume und in der übersichtlichsten zu beliebiger Prüfung der Richtigkeit geeignetsten Form besitzt.

\Article{21}
Ich gehe jetzt zu der Hauptaufgabe selbst über, welche für die Ellipsoidfläche eine ähnliche Methode fordert, wie für die Kugelfläche im Vorhergehenden gegeben ist. Die Auflösung dieser allerdings etwas verwickelten Aufgabe soll hier auf zwei ganz von einander verschiedenen Wegen abgeleitet werden. Da die eine Ableitung, mit welcher der Anfang gemacht werden wird, sich auf diejenige conforme Übertragung der Ellipsoidfläche auf die Kugelfläche gründet, deren Theorie in der ersten Abhandlung entwickelt ist, so kann die Auffindung dieser Auflösung wie die \textit{erste mittelbare} Benutzung dieser Theorie für die Zwecke der höhern Geodaesie betrachtet werden. (Vergl. Art. 11).

Es mögen demnach jetzt durch \(B+\frac12b\) und \(B-\frac12b\) die Breiten zweier Punkte auf der Ellipsoidfläche bezeichnet werden; ihr Längenunterschied durch \(l\); das zwischen ihnen enthaltene Stück einer geodaetischen Linie (und zwar hier nach beliebiger Einheit gemessen) durch \(r\); die Azimuthe der Linie am ersten und zweiten Endpunkte durch \(T+\frac12t\) und \(T-\frac12t\pm180^0\). Es handelt sich also darum, \(b,l\) und \(t\) aus \(r,B\) und \(T\) zu finden durch Formeln, welche den oben
